\documentclass[11pt,a4paper]{article}

% --- 核心宏包与中文支持 ---
\usepackage[utf8]{inputenc}
\usepackage{xeCJK}
\usepackage{amsmath, amssymb, amsfonts}
\usepackage{listings}
\usepackage{xcolor}
\usepackage{geometry}
\usepackage{tcolorbox}
\usepackage{booktabs}
\usepackage{tabularx}
\usepackage{setspace}
\usepackage{enumitem}
\usepackage[hidelinks]{hyperref}
\usepackage{bookmark}
\usepackage{fancyhdr}

% --- PDF 书签与元数据 ---
\hypersetup{
    colorlinks=false,
    pdfborder={0 0 0},
    bookmarks=true,
    bookmarksnumbered=true,
    bookmarksopen=true,
    bookmarksopenlevel=2,
    pdfcreator={XeLaTeX},
    pdftitle={01 - Variables & Data Types},
    pdfauthor={黄子扬},
}

% --- 页面美化设置 ---
\geometry{top=0.7in,bottom=0.8in,left=0.6in,right=0.6in}
\onehalfspacing
\tcbuselibrary{skins,breakable}
\setlength{\parindent}{0pt}
\setlength{\parskip}{0.6em}

% --- 代码样式配置 ---
\definecolor{mainblue}{RGB}{0,51,102}
\definecolor{examred}{RGB}{180,0,0}
\definecolor{codebg}{RGB}{248,248,248}
\definecolor{commentgreen}{RGB}{0,128,0}

\lstset{
    backgroundcolor=\color{codebg},
    basicstyle=\ttfamily\footnotesize,
    keywordstyle=\color{blue}\bfseries,
    commentstyle=\color{commentgreen},
    stringstyle=\color{orange},
    breaklines=true,
    showstringspaces=false,
    frame=leftline,
    xleftmargin=2em,
    numbers=left,
    numbersep=5pt,
    numberstyle=\tiny\color{gray},
    language=Python,
    morekeywords={None, True, False}
}

% --- 自定义教学框 ---
\newtcolorbox{concept}[1]{colback=blue!5, colframe=mainblue, fonttitle=\bfseries, title={{【Concept / 核心概念】#1}}, arc=3pt, breakable}
\newtcolorbox{warning}[1]{colback=red!5, colframe=examred, fonttitle=\bfseries, title={{【Warning / 避坑指南】#1}}, arc=3pt, breakable}
\newtcolorbox{examplebox}[1]{colback=green!2, colframe=green!40!black, fonttitle=\bfseries, title={{【Example / 实操案例】#1}}, arc=2pt, breakable}

% --- 页眉页脚 ---
\pagestyle{fancy}
\fancyhf{}
\fancyhead[L]{\small\bfseries 01 - Variables \& Data Types / 变量与数据类型}
\fancyhead[R]{\small\thepage}
\renewcommand{\headrulewidth}{0.5pt}

% --- 文档信息 ---
\title{\Huge \bfseries 01 - Variables \& Data Types}
\author{\Large \textbf{哈尔滨工业大学 \quad 黄子扬}}
\date{2026年6月8日}

\begin{document}

\maketitle
\tableofcontents
\newpage

% ======================================================================
\section{Variables / 变量基础}
% ======================================================================

\begin{concept}{What is a Variable? / 什么是变量？}
    In Python, variables are names that can be assigned a value and then used to refer to that value throughout your code. \\
    在 Python 中，变量是可以被赋值并在代码中反复引用的名称。

    Variables are fundamental to programming for two reasons:
    \begin{enumerate}[itemsep=0.3em]
        \item \textbf{Keep values accessible / 保持值可访问}：You can assign the result of a time-consuming operation to a variable so your program doesn't have to perform the operation each time you need the result.
        \item \textbf{Give values context / 赋予值上下文}：The number 28 could mean lots of different things. Giving it a name like \texttt{num\_students} makes the meaning clear.
    \end{enumerate}
\end{concept}

\subsection{The Assignment Operator / 赋值运算符}
\begin{concept}{How it works / 赋值逻辑}
    The \texttt{=} symbol is called the \textbf{assignment operator}. It takes the value on the \textbf{right} and assigns it to the name on the \textbf{left}. \\
    Python 使用 \texttt{=} 符号进行赋值。它将等号右侧的值"绑定"到左侧的名字上。
\end{concept}

\begin{examplebox}{Basic Assignment / 基础赋值}
\begin{lstlisting}
age = 25
print(age)      # Output: 25
\end{lstlisting}
\end{examplebox}

\begin{concept}{How print() Works / print() 的工作原理}
    \texttt{print()} is a function that displays output. Python looks for the variable name, finds its assigned value, and \textbf{replaces} the variable name with its value before calling the function. \\
    \texttt{print()} 是一个函数，用于显示输出。Python 先查找变量名对应的值，\textbf{用值替换}变量名后再调用函数。
\end{concept}

\subsection{Naming Rules \& Conventions / 命名规则与规范}

\begin{itemize}[itemsep=0.3em]
    \item \textbf{Valid Characters / 合法字符}：Uppercase and lowercase letters (A–Z, a–z), digits (0–9), and underscores (\_). Variable names can be as long or as short as you like.
    \item \textbf{Case Sensitivity / 大小写敏感}：\texttt{age} is not the same as \texttt{Age}. Using \texttt{print(Age)} when only \texttt{age} is defined will raise a \texttt{NameError}.
    \item \textbf{Style / 规范}：Common practice is to write variable names in \texttt{lower\_case\_with\_underscores} (e.g., \texttt{car\_length}).
\end{itemize}

\begin{warning}{Variable Name Rules / 变量命名禁忌}
    \begin{itemize}
        \item \textbf{No Digit Start / 严禁数字开头}：Names cannot begin with a digit. \texttt{3text = 4} is invalid.
        \item \textbf{Reserved Words / 保留字}：Do NOT name variables \texttt{print}, \texttt{for}, \texttt{in}, \texttt{list}, \texttt{str}, \texttt{int} — these are built-in functions you will override.
        \item \textbf{Meaningful Names / 命名要有意义}：Always use names that describe the data (e.g., \texttt{car\_length} instead of \texttt{L}).
        \item \textbf{Don't make names too long / 不要太长}：Find a balance between descriptiveness and length.
    \end{itemize}
\end{warning}

\begin{examplebox}{Valid and Invalid Names / 合法与非法命名}
\begin{lstlisting}
# Valid names / 合法命名
test1 = 1
_alpha = 2
L = 3
car_length = 2000

# Invalid names / 非法命名
# 3text = 4         # Error: cannot start with digit
# my-var = 5        # Error: hyphen is not allowed
\end{lstlisting}
\end{examplebox}

\subsection{Comments / 代码注释}

\begin{concept}{What are Comments? / 什么是注释？}
    Comments are lines of text that don't affect how a program runs. They document what code does or why the programmer made certain decisions. \\
    注释是不影响程序运行的文本行，用于记录代码用途或编程决策的原因。

    The most common way to write a comment is to begin a new line with the \texttt{\#} character. When Python runs your code, it ignores lines starting with \texttt{\#}.
\end{concept}

\begin{examplebox}{Comment Examples / 注释示例}
\begin{lstlisting}
# This variable stores the car length in cm
L = 2000
print(L)

H = 146   # This variable stores the car height in cm
print(H)
\end{lstlisting}
\end{examplebox}

% ======================================================================
\section{Data Types / 数据类型百科}
% ======================================================================

\begin{concept}{What is a Data Type? / 什么是数据类型？}
    The term \textbf{data type} refers to what kind of data a value represents. Python has several built-in data types. \\
    \textbf{数据类型}指一个值代表什么类型的数据。Python 有多种内置数据类型。
\end{concept}

\subsection{Primitive Types / 原始类型}

\begin{table}[h]
\centering
\renewcommand{\arraystretch}{1.5}
\begin{tabular}{@{}llll@{}}
\toprule
\textbf{Type / 类型} & \textbf{Keyword} & \textbf{Description / 描述} & \textbf{Example / 示例} \\ \midrule
Integer / 整数 & \texttt{int} & Whole numbers / 整数 & \texttt{100} \\
Float / 浮点数 & \texttt{float} & Decimal numbers / 小数 & \texttt{10.7} \\
Boolean / 布尔值 & \texttt{bool} & \texttt{True} or \texttt{False} (Capitalized!) & \texttt{True} \\
String / 字符串 & \texttt{str} & Text data / 引号内的文本 & \texttt{"Hello"} \\ \bottomrule
\end{tabular}
\end{table}

\subsection{The type() Function / type() 函数}

\begin{concept}{Determining Data Type / 确定数据类型}
    The \texttt{type()} function is used to determine the data type of a given value or variable. \\
    \texttt{type()} 函数用于判断任何值或变量的数据类型。
\end{concept}

\begin{examplebox}{Using type() / type() 实操}
\begin{lstlisting}
num = 100               # An integer variable
length = 10.7           # A floating point variable

print(type(num))        # <class 'int'>
print(type(length))     # <class 'float'>
\end{lstlisting}
\end{examplebox}

\subsection{Boolean Data Type / 布尔类型}

\begin{concept}{True or False / 真与假}
    The boolean data type (\texttt{bool}) can have only one of two values: \texttt{True} or \texttt{False} (both must be \textbf{capitalized}!). \\
    布尔型只有两个值：\texttt{True}（真）和 \texttt{False}（假），首字母\textbf{必须大写}！

    Booleans are fundamental for conditional expressions — we will explore them more in later chapters.
\end{concept}

\begin{examplebox}{Boolean Examples / 布尔值示例}
\begin{lstlisting}
is_finished = True
print(type(is_finished))    # <class 'bool'>

is_finished = False
print(type(is_finished))    # <class 'bool'>
\end{lstlisting}
\end{examplebox}

% ======================================================================
\section{Strings: Complete Guide / 字符串完全指南}
% ======================================================================

\subsection{Properties \& Creation / 属性与创建}

\begin{concept}{Three Key Properties / 三大核心属性}
    Strings in Python are identified as a contiguous set of characters represented in quotation marks. They have three important properties:
    \begin{enumerate}[itemsep=0.3em]
        \item \textbf{Characters / 字符}：Strings contain individual letters or symbols called characters.
        \item \textbf{Length / 长度}：Strings have a length, defined as the number of characters the string contains. Use \texttt{len()} to get it.
        \item \textbf{Sequence / 序列}：Characters in a string appear in a sequence — each character has a \textbf{numbered position} (index) in the string.
    \end{enumerate}
    字符串是引号内的连续字符集。三个核心属性：(1)由字符组成 (2)有长度 (3)字符按序列排列。
\end{concept}

\subsection{Quotes and Escape Characters / 引号与转义字符}

\begin{warning}{Quote Conflicts / 引号冲突}
    If you use double quotes inside a string delimited by double quotes, you'll get an error:
    \begin{lstlisting}
# text1 = "She said, "What time is it?""  # ERROR!
    \end{lstlisting}
    \textbf{Solutions / 解决方案}：
    \begin{enumerate}
        \item Use single quotes outside, double quotes inside:
        \texttt{text1 = 'She said, "What time is it?"'}
        \item Use the escape character \texttt{\textbackslash}:
        \texttt{text1 = "She said, \textbackslash"What time is it?\textbackslash""}
    \end{enumerate}
\end{warning}

\begin{examplebox}{Escape Character / 转义字符示例}
\begin{lstlisting}
text1 = 'She said, "What time is it?"'
text2 = "She said, \"What time is it?\""
print(text1)    # She said, "What time is it?"
print(text2)    # She said, "What time is it?"
\end{lstlisting}
\end{examplebox}

\subsection{Determine the Length of a String / 获取字符串长度}

\begin{examplebox}{Using len() / len() 函数}
    \texttt{len()} is a Python built-in function that returns the number of characters in a string (including spaces).
\begin{lstlisting}
text2 = "This is a test string"
print(len(text2))   # Output: 22
\end{lstlisting}
\end{examplebox}

\subsection{Multiline Strings / 多行字符串}

\begin{concept}{Two Ways to Create Multiline Strings / 两种多行字符串方式}
    To deal with long strings, you can break them up across multiple lines:
    \begin{enumerate}[itemsep=0.3em]
        \item \textbf{Backslash method / 反斜杠法}：Put a backslash (\texttt{\textbackslash}) at the end of each line except the last.
        \item \textbf{Triple quotes method / 三引号法}：Use triple quotes (\texttt{"""} or \texttt{'''}) as delimiters.
    \end{enumerate}
\end{concept}

\begin{examplebox}{Multiline String Examples / 多行字符串示例}
\begin{lstlisting}
# Method 1: Backslash / 反斜杠法
long_text = "This planet has--or rather had--a problem, which was \
this: most of the people living on it were unhappy for pretty much of the time."
print(long_text)

# Method 2: Triple quotes / 三引号法
long_text = """This planet has--or rather had--a problem, which was
this: most of the people living on it were unhappy for pretty much of the time."""
print(long_text)
\end{lstlisting}
\end{examplebox}

\subsection{Concatenation, Indexing, and Slicing / 拼接、索引与切片}

\subsubsection{Concatenation / 字符串拼接}
\begin{concept}{Joining Strings with + / 用 + 拼接字符串}
    You can combine, or concatenate, two strings using the \texttt{+} operator.
\end{concept}

\begin{examplebox}{Concatenation Example / 拼接示例}
\begin{lstlisting}
string1 = "Hello"
string2 = "World"
string3 = string1 + string2
print(string3)      # Output: HelloWorld (no space!)
\end{lstlisting}
\end{examplebox}

\subsubsection{Indexing / 索引}
\begin{concept}{The Rule of 0 / 从 0 开始计数}
    In Python, counting always starts at \textbf{zero (0)}. Each character in a string has a numbered position called an \textbf{index}. \\
    Python 计数永远从 0 开始。字符串中的每个字符都有一个称为索引的编号位置。
\end{concept}

\begin{examplebox}{Positive and Negative Indexing / 正索引与负索引}
\begin{lstlisting}
text = "Python BootCamp"

# Positive index / 正索引（从左到右，从 0 开始）
print(text[0])      # 'P' (第一个字符)
print(text[7])      # 'B'

# Negative index / 负索引（从右到左，从 -1 开始）
print(text[-1])     # 'p' (最后一个字符)
print(text[-5])     # 't'
\end{lstlisting}
\end{examplebox}

\subsubsection{Slicing / 切片}
\begin{concept}{Extracting Substrings / 提取子字符串}
    You can extract a portion of a string (called a \textbf{substring}) by inserting a colon between two index numbers inside square brackets: \texttt{[start:stop]}. \\
    \textbf{Critical Rule / 关键规则}：Inclusive on the left, \textbf{EXCLUSIVE} on the right. \\
    切片格式为 \texttt{[start:stop]}，左闭右开——包含 start，不包含 stop。
\end{concept}

\begin{examplebox}{Slicing Examples / 切片实战}
\begin{lstlisting}
text = "Python BootCamp"

# Basic slicing [start:stop]
print(text[0:6])    # 'Python' (indices 0,1,2,3,4,5 — NOT 6)
print(text[7:11])   # 'Boot'

# Omit start: assumes 0
print(text[:6])     # 'Python' (same as text[0:6])

# Omit stop: goes to the end
print(text[7:])     # 'BootCamp'

# Both omitted: the whole string
print(text[:])      # 'Python BootCamp'
\end{lstlisting}
\end{examplebox}

\subsection{Immutability / 不可变性}

\begin{warning}{Strings Are Immutable / 字符串不可修改}
    You cannot change characters once a string is created. Strings are \textbf{immutable}. \\
    字符串一旦创建就不可修改。试图执行 \texttt{text[6] = '\_'} 会抛出 \texttt{TypeError}：\texttt{'str' object does not support item assignment}.

    \textbf{Solution / 解决方法}：Most string methods return a \textbf{modified copy} — you must reassign to keep the result.
\end{warning}

\begin{examplebox}{Immutability Demonstration / 不可变性演示}
\begin{lstlisting}
name = "Emlyon"
name.upper()            # Returns "EMLYON" but does NOT change 'name'!
print(name)             # Still "Emlyon"

name = name.upper()     # Correct: reassign the result
print(name)             # Now "EMLYON"
\end{lstlisting}
\end{examplebox}

\subsection{String Methods / 字符串方法大全}

\begin{concept}{What are String Methods? / 什么是字符串方法？}
    Strings come bundled with special functions called \textbf{string methods} that you can use to manipulate strings. The dot (\texttt{.}) tells Python that what follows is the name of a method. \\
    字符串附带特殊的函数称为字符串方法，点号 (\texttt{.}) 表示接下来是方法名。
\end{concept}

\begin{examplebox}{Common String Methods / 常用字符串方法}
\begin{lstlisting}
# 1. Case Conversion / 大小写转换
print("Python BootCamp".lower())        # 'python bootcamp'
print("Python BootCamp".upper())        # 'PYTHON BOOTCAMP'

# 2. Removing Whitespace / 去除首尾空白
print(" Python BootCamp".strip())       # 'Python BootCamp'

# 3. Counting Substrings / 统计子串出现次数
text = "Python BootCamp is the best Python program ever!"
print(text.count("Python"))             # Output: 2

# 4. Finding Substring Position / 查找子串位置
print(text.find("Python"))              # Output: 0 (first occurrence)
print(text.find("python"))              # Output: -1 (not found!)

# 5. Check Start and End / 检查首尾
print("Python BootCamp".startswith('P'))  # True
print("Python BootCamp".endswith('p'))    # False

# 6. Replace / 替换
print("Hello World".replace("World", "Python"))  # 'Hello Python'
\end{lstlisting}
\end{examplebox}

\begin{warning}{find() Returns -1 When Not Found / find() 找不到返回 -1}
    The \texttt{find()} method returns the index of the first occurrence, or \textbf{-1} if the substring is not found. This is different from indexing with \texttt{[]}, which would raise an error. \\
    \texttt{find()} 找不到时返回 \texttt{-1}（不是报错），与直接用 \texttt{[]} 索引不同。
\end{warning}

% ======================================================================
\section{The None Data Type / 特殊类型 None}
% ======================================================================

\begin{concept}{NoneType / 空值类型}
    The \texttt{None} keyword is used to define a null value, or no value at all. \\
    \texttt{None} 用于表示"空值"或"无值"。

    \textbf{Critical Distinctions / 关键区分}：
    \begin{itemize}
        \item \texttt{None} is NOT the same as \texttt{0}
        \item \texttt{None} is NOT the same as \texttt{False}
        \item \texttt{None} is NOT the same as an empty string \texttt{""}
        \item \texttt{None} is a data type of its own (\texttt{NoneType})
        \item Only \texttt{None} can be \texttt{None}
    \end{itemize}
\end{concept}

\begin{examplebox}{None Example / None 示例}
\begin{lstlisting}
my_var = None
print(my_var)       # None
print(type(my_var)) # <class 'NoneType'>
\end{lstlisting}
\end{examplebox}

% ======================================================================
\section{Operators / 运算符}
% ======================================================================

\subsection{Arithmetic Operators / 算术运算符}

\begin{table}[h]
\centering
\renewcommand{\arraystretch}{1.5}
\begin{tabular}{@{}lll@{}}
\toprule
\textbf{Operator} & \textbf{Operation / 操作} & \textbf{Example / 示例} \\ \midrule
\texttt{+}  & Addition / 加法          & \texttt{2 + 2 = 4} \\
\texttt{-}  & Subtraction / 减法       & \texttt{5 - 2 = 3} \\
\texttt{*}  & Multiplication / 乘法    & \texttt{3 * 2 = 6} \\
\texttt{/}  & Division / 标准除法      & \texttt{5 / 2 = 2.5} \\
\texttt{//} & Floor Division / 整除    & \texttt{5 // 2 = 2} \\
\texttt{**} & Exponent / 幂运算        & \texttt{3 ** 2 = 9} \\ \bottomrule
\end{tabular}
\end{table}

\begin{warning}{/ vs // 的区别}
    In Python 3, \texttt{/} denotes \textbf{standard division} (always returns float), while \texttt{//} denotes \textbf{floor division} (integer division, discards the remainder). \\
    Python 3 中，\texttt{/} 是标准除法（返回 float），\texttt{//} 是整除（丢弃余数）。
\end{warning}

\begin{examplebox}{Arithmetic Operators Demo / 算术运算符演示}
\begin{lstlisting}
print(2 + 2)        # 4
print(3 * 2)        # 6
print(3 ** 2)       # 9
print(5 / 2)        # 2.5 (standard division)
print(5 // 2)       # 2 (floor division)
\end{lstlisting}
\end{examplebox}

\subsection{Arithmetic on Strings / 字符串上的算术运算}

\begin{concept}{Operators work differently on strings / 运算符对字符串有特殊行为}
    Some arithmetic operators also work on \texttt{str} type:
    \begin{itemize}
        \item \texttt{+} on strings: concatenation / 拼接
        \item \texttt{*} between string and int: repetition / 重复
    \end{itemize}
\end{concept}

\begin{examplebox}{String Arithmetic / 字符串运算}
\begin{lstlisting}
print('2' + '4')        # '24' (concatenation, not addition!)
print(2 * 'S')          # 'SS' (repetition)
# print('2' * '4')      # Error! Cannot multiply string by string
\end{lstlisting}
\end{examplebox}

\subsection{Comparison Operators / 比较运算符}

\begin{concept}{Comparison Produces Boolean / 比较结果总是布尔值}
    The result of a comparison is always a boolean value (\texttt{True} or \texttt{False}). \\
    比较运算的结果永远是布尔值（\texttt{True} 或 \texttt{False}）。
\end{concept}

\begin{table}[h]
\centering
\renewcommand{\arraystretch}{1.5}
\begin{tabular}{@{}lll@{}}
\toprule
\textbf{Operator} & \textbf{Description / 描述} & \textbf{Example / 示例} \\ \midrule
\texttt{==} & Equal / 等于               & \texttt{1 == 1} $\rightarrow$ \texttt{True} \\
\texttt{!=} & Not equal / 不等于          & \texttt{2.3 != 2.313} $\rightarrow$ \texttt{True} \\
\texttt{<}  & Less than / 小于             & \texttt{2.5 < 9} $\rightarrow$ \texttt{True} \\
\texttt{>}  & Greater than / 大于          & \texttt{2.4 > 2.399} $\rightarrow$ \texttt{True} \\
\texttt{<=} & Less than or equal / 小于等于  & \texttt{3 <= 3} $\rightarrow$ \texttt{True} \\
\texttt{>=} & Greater than or equal / 大于等于 & \texttt{2.4 >= 2.399} $\rightarrow$ \texttt{True} \\ \bottomrule
\end{tabular}
\end{table}

\begin{examplebox}{Comparison Examples / 比较运算示例}
\begin{lstlisting}
print(2.5 < 3)          # True
print(2 == 2)           # True
print(2 != 3)           # True
print(5 >= 5)           # True
print(10 < 3)           # False
\end{lstlisting}
\end{examplebox}

\subsection{Logical Operators / 逻辑运算符}

\begin{concept}{Combining Boolean Values / 组合布尔值}
    Logical operators take boolean values as input and produce either \texttt{True} or \texttt{False} as output. \\
    逻辑运算符接收布尔值并返回布尔值。
\end{concept}

\begin{table}[h]
\centering
\renewcommand{\arraystretch}{1.5}
\begin{tabularx}{\textwidth}{@{}l>{\raggedright\arraybackslash}X>{\raggedright\arraybackslash}X@{}}
\toprule
\textbf{Operator} & \textbf{Description / 描述} & \textbf{Rule / 规则} \\ \midrule
\texttt{x or y}  & If either is True $\rightarrow$ True; otherwise False. / 任一为真则真 & Returns \texttt{True} if at least one is \texttt{True} \\
\texttt{x and y} & If both are True $\rightarrow$ True; otherwise False. / 两者为真才真 & Returns \texttt{True} only if \textbf{both} are \texttt{True} \\
\texttt{not x}   & Reverses the boolean value. / 取反 & Returns \texttt{False} if \texttt{x} is \texttt{True}; \texttt{True} if \texttt{x} is \texttt{False} \\ \bottomrule
\end{tabularx}
\end{table}

\begin{examplebox}{Logical Operators Demo / 逻辑运算符演示}
\begin{lstlisting}
# or: True if either is True
print((2 == 3) or (4 < 5))      # True (4 < 5 is True)

# and: True only if BOTH are True
print((2 == 3) and (4 < 5))     # False (2 == 3 is False)

# not: Reverses the boolean
print(not (2 == 3))             # True (reverses False to True)
print(not True)                 # False
\end{lstlisting}
\end{examplebox}

% ======================================================================
\section{Type Casting / 类型转换}
% ======================================================================

\begin{concept}{Converting Between Types / 类型之间的转换}
    If you want to change one data type to another, use the following functions: \\
    使用以下函数进行类型转换：\texttt{int()}, \texttt{str()}, \texttt{float()}, \texttt{bool()}.

    These four data types (\texttt{int}, \texttt{float}, \texttt{str}, \texttt{bool}) are called \textbf{Primitive Types}.
\end{concept}

\begin{examplebox}{Type Casting Examples / 类型转换示例}
\begin{lstlisting}
print(int(9.9))         # 9 (truncates, does NOT round!)
print(str(5))           # '5'
print(bool(1))          # True
# print(int("Yes!"))    # ValueError: invalid literal for int()
\end{lstlisting}
\end{examplebox}

\begin{warning}{String + Int = Error! / 字符串 + 整数 = 报错！}
    You MUST cast \texttt{str} to \texttt{int} before doing arithmetic operations:
\begin{lstlisting}
# Wrong / 错误
a = 34
b = "10"
# print(a + b)        # TypeError! Cannot add int and str

# Correct / 正确
a = 34
b = "10"
print(a + int(b))       # 44
\end{lstlisting}
\end{warning}

\begin{concept}{Truthiness / 真值判定规则}
    Almost any value is evaluated to \texttt{True} if it has some sort of content:
    \begin{itemize}[itemsep=0.3em]
        \item Any string is \texttt{True}, \textbf{except} empty strings (\texttt{""}).
        \item Any number is \texttt{True}, \textbf{except} \texttt{0}.
        \item \texttt{bool(0)} $\rightarrow$ \texttt{False}; \texttt{bool("")} $\rightarrow$ \texttt{False}
    \end{itemize}
    \textbf{Falsy values / 假值}：\texttt{0}, \texttt{0.0}, \texttt{""}, \texttt{None}, \texttt{False}, empty containers.
\end{concept}

% ======================================================================
\section{String Formatting / 字符串格式化输出（四种方法）}
% ======================================================================

\begin{concept}{The Goal / 目标}
    Suppose you have \texttt{name = "Doctor Strange"}, \texttt{heads = 2}, \texttt{arms = 3}. \\
    You want to display: \textbf{"Doctor Strange has 2 heads and 3 arms."}

    There are \textbf{four} ways to achieve this in Python.
\end{concept}

\subsection{Method 1: Commas in print() / 用逗号分隔}

\begin{examplebox}{Comma Method / 逗号法}
\begin{lstlisting}
name = "Doctor Strange"
heads = 2
arms = 3
print(name, "has", str(heads), "heads and", str(arms), "arms.")
# Output: Doctor Strange has 2 heads and 3 arms.
\end{lstlisting}
    Note: \texttt{print()} automatically inserts spaces between comma-separated items. You must convert numbers to strings with \texttt{str()}.
\end{examplebox}

\subsection{Method 2: Concatenation with + / 字符串拼接法}

\begin{examplebox}{Concatenation Method / 拼接法}
\begin{lstlisting}
print(name + " has " + str(heads) + " heads and " + str(arms) + " arms.")
# Output: Doctor Strange has 2 heads and 3 arms.
\end{lstlisting}
    Note: Tedious! Must manually add spaces and convert numbers to strings.
\end{examplebox}

\subsection{Method 3: f-strings (Best Practice, Python 3.6+) / f-字符串（推荐）}

\begin{concept}{f-strings / f-字符串格式化}
    Two important things to notice:
    \begin{enumerate}
        \item The string literal starts with the letter \texttt{f} before the opening quotation mark.
        \item Variable names surrounded by curly braces \texttt{\{} and \texttt{\}} are replaced with their corresponding values \textbf{without} needing \texttt{str()}.
    \end{enumerate}
    f-strings are only available in \textbf{Python version 3.6 and above}.
\end{concept}

\begin{examplebox}{f-string Example / f-字符串示例}
\begin{lstlisting}
print(f"{name} has {heads} heads and {arms} arms.")
# Output: Doctor Strange has 2 heads and 3 arms.
\end{lstlisting}
\end{examplebox}

\subsection{Method 4: .format() Method / .format() 方法}

\begin{examplebox}{.format() Method / .format() 方法示例}
\begin{lstlisting}
print("{} has {} heads and {} arms.".format(name, heads, arms))
# Output: Doctor Strange has 2 heads and 3 arms.
\end{lstlisting}
    Curly braces \texttt{\{\}} act as placeholders, filled by \texttt{.format()} arguments in order.
\end{examplebox}

\subsection{Method 5 (Legacy): \% Operator / \% 运算符（旧式）}

\begin{examplebox}{\% Operator / \% 运算符示例}
\begin{lstlisting}
print("%s has %s heads and %s arms." % (name, heads, arms))
# Output: Doctor Strange has 2 heads and 3 arms.
\end{lstlisting}
    \texttt{\%s} is a placeholder for string; values are provided in parentheses after \texttt{\%}.
\end{examplebox}

% ======================================================================
\section{User Input / 用户输入}
% ======================================================================

\begin{concept}{The input() Function / input() 函数}
    The \texttt{input()} function gets input from the user. It:
    \begin{itemize}[itemsep=0.3em]
        \item Displays a \textbf{prompt} (a string that tells the user what to enter).
        \item Waits for the user to type something and press Enter.
        \item \textbf{Always returns a string} — even if the user types a number!
    \end{itemize}
\end{concept}

\begin{examplebox}{input() Examples / 用户输入示例}
\begin{lstlisting}
# Basic usage
name = input()
print(name)

# With a prompt / 带提示信息的输入
age = input("Please enter your age: ")
print(age)              # age is a STRING, not a number!
print(type(age))        # <class 'str'>

# Convert to number if needed / 如需数字，必须转换
age = int(input("Please enter your age: "))
print(age + 1)          # Now it works as a number
\end{lstlisting}
\end{examplebox}

\begin{warning}{input() Always Returns a String / input() 永远返回字符串}
    Even if the user types \texttt{42}, the result is the string \texttt{"42"}, NOT the integer \texttt{42}. Use \texttt{int()} or \texttt{float()} to convert. \\
    即使用户输入 42，返回的也是字符串 \texttt{"42"}，需要用 \texttt{int()} 或 \texttt{float()} 转换后才能做数学运算。
\end{warning}

% ======================================================================
\section{Quick Reference / 速查表}
% ======================================================================

\begin{table}[h]
\centering
\small
\renewcommand{\arraystretch}{1.3}
\begin{tabular}{@{}ll@{}}
\toprule
\textbf{Operation / 操作} & \textbf{Syntax / 语法} \\ \midrule
赋值 & \texttt{x = 10} \\
检查类型 & \texttt{type(x)} \\
获取长度 & \texttt{len(s)} \\
索引（正） & \texttt{s[0]} \\
索引（负） & \texttt{s[-1]} \\
切片 & \texttt{s[start:stop]} \\
小写 & \texttt{s.lower()} \\
大写 & \texttt{s.upper()} \\
去空白 & \texttt{s.strip()} \\
计数 & \texttt{s.count(sub)} \\
查找 & \texttt{s.find(sub)} \\
开头判断 & \texttt{s.startswith(prefix)} \\
结尾判断 & \texttt{s.endswith(suffix)} \\
替换 & \texttt{s.replace(old, new)} \\
转整数 & \texttt{int(x)} \\
转浮点 & \texttt{float(x)} \\
转字符串 & \texttt{str(x)} \\
转布尔 & \texttt{bool(x)} \\
f-字符串 & \texttt{f"\{var\}"} \\
用户输入 & \texttt{input(prompt)} \\
整除 & \texttt{//} \\
幂运算 & \texttt{**} \\
逻辑与 & \texttt{and} \\
逻辑或 & \texttt{or} \\
逻辑非 & \texttt{not} \\
等于 & \texttt{==} \\
不等于 & \texttt{!=} \\ \bottomrule
\end{tabular}
\end{table}

\end{document}
